\section{Projective morphisms}
\label{section-projective}

\noindent
We will use the definition of a projective morphism from \cite{EGA}.
The version of the definition with the ``H'' is the one
from \cite{H}. The resulting definitions are different. Both are useful.

\begin{definition}
\label{definition-projective}
Let $f : X \to S$ be a morphism of schemes.
\begin{enumerate}
\item We say $f$ is {\it projective} if $X$ is isomorphic as
an $S$-scheme to a closed subscheme of a projective
bundle $\mathbf{P}(\mathcal{E})$
for some quasi-coherent, finite type $\mathcal{O}_S$-module $\mathcal{E}$.
\item We say $f$ is {\it H-projective} if there exists an integer $n$ and
a closed immersion $X \to \mathbf{P}^n_S$ over $S$.
\item We say $f$ is {\it locally projective} if there exists an open
covering $S = \bigcup U_i$ such that each $f^{-1}(U_i) \to U_i$ is
projective.
\end{enumerate}
\end{definition}

\noindent
As expected, a projective morphism is quasi-projective, see
Lemma \ref{lemma-projective-quasi-projective}.
Conversely, quasi-projective morphisms are often compositions
of open immersions and projective morphisms, see
Lemma \ref{lemma-quasi-projective-open-projective}.
For an overview of properties of projective morphisms over
a quasi-projective base, see
More on Morphisms, Section \ref{more-morphisms-section-projective}.

\begin{example}
\label{example-projective}
Let $S$ be a scheme. Let $\mathcal{A}$ be a quasi-coherent graded
$\mathcal{O}_S$-algebra generated by $\mathcal{A}_1$ over $\mathcal{A}_0$.
Assume furthermore that $\mathcal{A}_1$ is of finite type over
$\mathcal{O}_S$. Set $X = \underline{\text{Proj}}_S(\mathcal{A})$.
In this case $X \to S$ is projective. Namely,
the morphism associated to the graded $\mathcal{O}_S$-algebra map
$$
\text{Sym}_{\mathcal{O}_X}^*(\mathcal{A}_1)
\longrightarrow
\mathcal{A}
$$
is a closed immersion, see
Constructions, Lemma
\ref{constructions-lemma-surjective-generated-degree-1-map-relative-proj}.
\end{example}
 
\begin{lemma}
\label{lemma-H-projective}
An H-projective morphism is H-quasi-projective.
An H-projective morphism is projective.
\end{lemma}

\begin{proof}
The first statement is immediate from the definitions.
The second holds as $\mathbf{P}^n_S$ is a projective bundle over $S$, see
Constructions, Lemma \ref{constructions-lemma-projective-space-bundle}.
\end{proof}

\begin{lemma}
\label{lemma-characterize-locally-projective}
Let $f : X \to S$ be a morphism of schemes.
The following are equivalent:
\begin{enumerate}
\item The morphism $f$ is locally projective.
\item There exists an open covering $S = \bigcup U_i$ such
that each $f^{-1}(U_i) \to U_i$ is H-projective.
\end{enumerate}
\end{lemma}

\begin{proof}
By Lemma \ref{lemma-H-projective} we see that (2) implies (1). Assume (1).
For every point $s \in S$ we can find $\Spec(R) = U \subset S$
an affine open neighbourhood of $s$ such that $X_U$ is isomorphic to a
closed subscheme of $\mathbf{P}(\mathcal{E})$ for some finite type,
quasi-coherent sheaf of $\mathcal{O}_U$-modules $\mathcal{E}$.
Write $\mathcal{E} = \widetilde{M}$ for some finite type
$R$-module $M$ (see
Properties, Lemma \ref{properties-lemma-finite-type-module}).
Choose generators $x_0, \ldots, x_n \in M$ of $M$ as an $R$-module.
Consider the surjective graded $R$-algebra map
$$
R[X_0, \ldots, X_n] \longrightarrow \text{Sym}_R(M).
$$
According to
Constructions, Lemma \ref{constructions-lemma-surjective-graded-rings-map-proj}
the corresponding morphism
$$
\mathbf{P}(\mathcal{E}) \to \mathbf{P}^n_R
$$
is a closed immersion. Hence we conclude that $f^{-1}(U)$ is isomorphic
to a closed subscheme of $\mathbf{P}^n_U$ (as a scheme over $U$).
In other words: (2) holds.
\end{proof}

\begin{lemma}
\label{lemma-locally-projective-proper}
A locally projective morphism is proper.
\end{lemma}

\begin{proof}
Let $f : X \to S$ be locally projective.
In order to show that $f$ is proper we may work locally on the
base, see Lemma \ref{lemma-proper-local-on-the-base}.
Hence, by Lemma \ref{lemma-characterize-locally-projective}
above we may assume there exists a closed immersion $X \to \mathbf{P}^n_S$.
By Lemmas \ref{lemma-composition-proper}
and \ref{lemma-closed-immersion-proper} it suffices to prove that
$\mathbf{P}^n_S \to S$ is proper. Since
$\mathbf{P}^n_S \to S$ is the base change of
$\mathbf{P}^n_{\mathbf{Z}} \to \Spec(\mathbf{Z})$ it suffices
to show that $\mathbf{P}^n_{\mathbf{Z}} \to \Spec(\mathbf{Z})$
is proper, see Lemma \ref{lemma-base-change-proper}.
By Constructions, Lemma \ref{constructions-lemma-proj-separated} the scheme
$\mathbf{P}^n_{\mathbf{Z}}$ is separated.
By Constructions, Lemma \ref{constructions-lemma-proj-quasi-compact} the scheme
$\mathbf{P}^n_{\mathbf{Z}}$ is quasi-compact.
It is clear that $\mathbf{P}^n_{\mathbf{Z}} \to \Spec(\mathbf{Z})$
is locally of finite type since $\mathbf{P}^n_{\mathbf{Z}}$ is
covered by the affine opens $D_{+}(X_i)$ each of which is the
spectrum of the finite type $\mathbf{Z}$-algebra
$$
\mathbf{Z}[X_0/X_i, \ldots, X_n/X_i].
$$
Finally, we have to show that
$\mathbf{P}^n_{\mathbf{Z}} \to \Spec(\mathbf{Z})$
is universally closed. This follows from
Constructions, Lemma \ref{constructions-lemma-proj-valuative-criterion}
and the valuative criterion (see Schemes,
Proposition \ref{schemes-proposition-characterize-universally-closed}).
\end{proof}

\begin{lemma}
\label{lemma-proper-ample-locally-projective}
Let $f : X \to S$ be a proper morphism of schemes. If there exists
an $f$-ample invertible sheaf on $X$, then $f$ is locally projective.
\end{lemma}

\begin{proof}
If there exists an $f$-ample invertible sheaf, then we can locally
on $S$ find an immersion $i : X \to \mathbf{P}^n_S$, see
Lemma \ref{lemma-finite-type-over-affine-ample-very-ample}. Since $X \to S$
is proper the morphism $i$ is a closed immersion, see
Lemma \ref{lemma-image-proper-scheme-closed}.
\end{proof}

\begin{lemma}
\label{lemma-H-projective-composition}
A composition of H-projective morphisms is H-projective.
\end{lemma}

\begin{proof}
Suppose $X \to Y$ and $Y \to Z$ are H-projective.
Then there exist closed immersions $X \to \mathbf{P}^n_Y$
over $Y$, and $Y \to \mathbf{P}^m_Z$ over $Z$.
Consider the following diagram
$$
\xymatrix{
X \ar[r] \ar[d] &
\mathbf{P}^n_Y \ar[r] \ar[dl] &
\mathbf{P}^n_{\mathbf{P}^m_Z} \ar[dl] \ar@{=}[r] &
\mathbf{P}^n_Z \times_Z \mathbf{P}^m_Z \ar[r] &
\mathbf{P}^{nm + n + m}_Z \ar[ddllll] \\
Y \ar[r] \ar[d] & \mathbf{P}^m_Z \ar[dl] & \\
Z & &
}
$$
Here the rightmost top horizontal arrow is the Segre embedding,
see Constructions, Lemma \ref{constructions-lemma-segre-embedding}.
The diagram identifies
$X$ as a closed subscheme of $\mathbf{P}^{nm + n + m}_Z$ as desired.
\end{proof}

\begin{lemma}
\label{lemma-H-projective-base-change}
A base change of a H-projective morphism is H-projective.
\end{lemma}

\begin{proof}
This is true because the base change of projective space
over a scheme is projective space, and the fact that the base
change of a closed immersion is a closed immersion, see
Schemes, Lemma \ref{schemes-lemma-base-change-immersion}.
\end{proof}

\begin{lemma}
\label{lemma-base-change-projective}
A base change of a (locally) projective morphism is (locally) projective.
\end{lemma}

\begin{proof}
This is true because the base change of a projective bundle
over a scheme is a projective bundle, the pullback of
a finite type $\mathcal{O}$-module is of finite type
(Modules, Lemma \ref{modules-lemma-pullback-finite-type})
and the fact that the base
change of a closed immersion is a closed immersion, see
Schemes, Lemma \ref{schemes-lemma-base-change-immersion}.
Some details omitted.
\end{proof}

\begin{lemma}
\label{lemma-projective-quasi-projective}
A projective morphism is quasi-projective.
\end{lemma}

\begin{proof}
Let $f : X \to S$ be a projective morphism. Choose a closed immersion
$i : X \to \mathbf{P}(\mathcal{E})$ where $\mathcal{E}$ is a quasi-coherent,
finite type $\mathcal{O}_S$-module. Then
$\mathcal{L} = i^*\mathcal{O}_{\mathbf{P}(\mathcal{E})}(1)$ is $f$-very ample.
Since $f$ is proper (Lemma \ref{lemma-locally-projective-proper})
it is quasi-compact. Hence Lemma \ref{lemma-ample-very-ample} implies
that $\mathcal{L}$ is $f$-ample. Since $f$ is proper it is of finite type.
Thus we've checked all the defining properties of quasi-projective
holds and we win.
\end{proof}

\begin{lemma}
\label{lemma-H-quasi-projective-open-H-projective}
Let $f : X \to S$ be a H-quasi-projective morphism.
Then $f$ factors as $X \to X' \to S$ where $X \to X'$ is an
open immersion and $X' \to S$ is H-projective.
\end{lemma}

\begin{proof}
By definition we can factor $f$ as a quasi-compact immersion
$i : X \to \mathbf{P}^n_S$ followed by the projection $\mathbf{P}^n_S \to S$.
By Lemma \ref{lemma-quasi-compact-immersion} there exists a closed
subscheme $X' \subset \mathbf{P}^n_S$ such that $i$ factors through
an open immersion $X \to X'$. The lemma follows.
\end{proof}

\begin{lemma}
\label{lemma-quasi-projective-open-projective}
Let $f : X \to S$ be a quasi-projective morphism with $S$ quasi-compact
and quasi-separated. Then $f$ factors as $X \to X' \to S$ where $X \to X'$
is an open immersion and $X' \to S$ is projective.
\end{lemma}

\begin{proof}
Let $\mathcal{L}$ be $f$-ample. Since $f$ is of finite type and $S$ is
quasi-compact $\mathcal{L}^{\otimes n}$ is $f$-very ample for some $n > 0$, see
Lemma \ref{lemma-finite-type-ample-very-ample}.
Replace $\mathcal{L}$ by $\mathcal{L}^{\otimes n}$.
Write $\mathcal{F} = f_*\mathcal{L}$. This is a quasi-coherent
$\mathcal{O}_S$-module by 
Schemes, Lemma \ref{schemes-lemma-push-forward-quasi-coherent}
(quasi-projective morphisms are quasi-compact and
separated, see Lemma \ref{lemma-quasi-projective-properties}). By
Properties, Lemma \ref{properties-lemma-directed-colimit-finite-presentation}
we can find a directed set $I$ and a system of
finite type quasi-coherent $\mathcal{O}_S$-modules $\mathcal{E}_i$
over $I$ such that $\mathcal{F} = \colim \mathcal{E}_i$.
Consider the compositions
$\psi_i : f^*\mathcal{E}_i \to f^*\mathcal{F} \to \mathcal{L}$.
Choose a finite affine open covering $S = \bigcup_{j = 1, \ldots, m} V_j$.
For each $j$ we can choose sections
$$
s_{j, 0}, \ldots, s_{j, n_j} \in
\Gamma(f^{-1}(V_j), \mathcal{L}) = f_*\mathcal{L}(V_j) = \mathcal{F}(V_j)
$$
which generate $\mathcal{L}$ over $f^{-1}V_j$ and define an immersion
$$
f^{-1}V_j \longrightarrow \mathbf{P}^{n_j}_{V_j},
$$
see Lemma \ref{lemma-very-ample-finite-type-over-affine}.
Choose $i$ such that there exist sections $e_{j, t} \in \mathcal{E}_i(V_j)$
mapping to $s_{j, t}$ in $\mathcal{F}$ for all $j = 1, \ldots, m$ and
$t = 1, \ldots, n_j$. Then the map $\psi_i$ is surjective
as the sections $f^*e_{j, t}$ have the same image as the sections $s_{j, t}$
which generate $\mathcal{L}|_{f^{-1}V_j}$. Whence we obtain a morphism
$$
r_{\mathcal{L}, \psi_i} : X \longrightarrow \mathbf{P}(\mathcal{E}_i)
$$
over $S$ such that over $V_j$ we have a factorization
$$
f^{-1}V_j \to \mathbf{P}(\mathcal{E}_i)|_{V_j} \to \mathbf{P}^{n_j}_{V_j}
$$
of the immersion given above. It follows that $r_{\mathcal{L}, \psi_i}|_{V_j}$
is an immersion, see Lemma \ref{lemma-immersion-permanence}.
Since $S = \bigcup V_j$ we conclude that $r_{\mathcal{L}, \psi_i}$
is an immersion.
Note that $r_{\mathcal{L}, \psi_i}$ is quasi-compact as
$X \to S$ is quasi-compact and $\mathbf{P}(\mathcal{E}_i) \to S$ is separated
(see Schemes, Lemma \ref{schemes-lemma-quasi-compact-permanence}).
By Lemma \ref{lemma-quasi-compact-immersion} there exists a closed
subscheme $X' \subset \mathbf{P}(\mathcal{E}_i)$ such that $i$ factors
through an open immersion $X \to X'$. Then $X' \to S$ is projective
by definition and we win.
\end{proof}

\begin{lemma}
\label{lemma-projective-is-quasi-projective-proper}
Let $S$ be a quasi-compact and quasi-separated scheme.
Let $f : X \to S$ be a morphism of schemes. Then
\begin{enumerate}
\item $f$ is projective if and only if $f$ is quasi-projective and proper, and
\item $f$ is H-projective if and only if $f$ is H-quasi-projective and proper.
\end{enumerate}
\end{lemma}

\begin{proof}
If $f$ is projective, then $f$ is quasi-projective by
Lemma \ref{lemma-projective-quasi-projective} and proper by
Lemma \ref{lemma-locally-projective-proper}. Conversely, if
$X \to S$ is quasi-projective and proper, then we can choose
an open immersion $X \to X'$ with $X' \to S$ projective by
Lemma \ref{lemma-quasi-projective-open-projective}.
Since $X \to S$ is proper, we see that $X$ is closed in $X'$
(Lemma \ref{lemma-image-proper-scheme-closed}), i.e.,
$X \to X'$ is a (open and) closed immersion.
Since $X'$ is isomorphic to a closed subscheme of a projective
bundle over $S$ (Definition \ref{definition-projective})
we see that the same thing is true for $X$, i.e.,
$X \to S$ is a projective morphism. This proves (1).
The proof of (2) is the same, except it uses
Lemmas \ref{lemma-H-projective} and
\ref{lemma-H-quasi-projective-open-H-projective}.
\end{proof}

\begin{lemma}
\label{lemma-composition-projective}
Let $f : X \to Y$ and $g : Y \to S$ be morphisms of schemes.
If $S$ is quasi-compact and quasi-separated and $f$ and $g$ are projective,
then $g \circ f$ is projective.
\end{lemma}

\begin{proof}
By Lemmas \ref{lemma-projective-quasi-projective} and
\ref{lemma-locally-projective-proper}
we see that $f$ and $g$ are quasi-projective and proper.
By Lemmas \ref{lemma-composition-proper} and
\ref{lemma-composition-quasi-projective}
we see that $g \circ f$ is proper and quasi-projective.
Thus $g \circ f$ is projective by
Lemma \ref{lemma-projective-is-quasi-projective-proper}.
\end{proof}

\begin{lemma}
\label{lemma-projective-permanence}
Let $g : Y \to S$ and $f : X \to Y$ be morphisms of schemes.
If $g \circ f$ is projective and $g$ is separated,
then $f$ is projective.
\end{lemma}

\begin{proof}
Choose a closed immersion $X \to \mathbf{P}(\mathcal{E})$ where $\mathcal{E}$
is a quasi-coherent, finite type $\mathcal{O}_S$-module. Then
we get a morphism $X \to \mathbf{P}(\mathcal{E}) \times_S Y$.
This morphism is a closed immersion because it is the composition
$$
X \to X \times_S Y \to \mathbf{P}(\mathcal{E}) \times_S Y
$$
where the first morphism is a closed immersion by
Schemes, Lemma \ref{schemes-lemma-semi-diagonal}
(and the fact that $g$ is separated) and
the second as the base change of a closed immersion.
Finally, the fibre product $\mathbf{P}(\mathcal{E}) \times_S Y$
is isomorphic to $\mathbf{P}(g^*\mathcal{E})$ and pullback
preserves quasi-coherent, finite type modules.
\end{proof}

\begin{lemma}
\label{lemma-projective-over-quasi-projective-is-H-projective}
Let $S$ be a scheme which admits an ample invertible sheaf. Then
\begin{enumerate}
\item any projective morphism $X \to S$ is H-projective, and
\item any quasi-projective morphism $X \to S$ is H-quasi-projective.
\end{enumerate}
\end{lemma}

\begin{proof}
The assumptions on $S$ imply that $S$ is quasi-compact and separated, see
Properties, Definition \ref{properties-definition-ample} and
Lemma \ref{properties-lemma-ample-immersion-into-proj}
and Constructions, Lemma \ref{constructions-lemma-proj-separated}.
Hence Lemma \ref{lemma-quasi-projective-open-projective}
applies and we see that (1) implies (2).
Let $\mathcal{E}$ be a finite type quasi-coherent $\mathcal{O}_S$-module.
By our definition of projective morphisms it suffices to show that
$\mathbf{P}(\mathcal{E}) \to S$ is H-projective.
If $\mathcal{E}$ is generated by finitely many global sections,
then the corresponding surjection $\mathcal{O}_S^{\oplus n} \to \mathcal{E}$
induces a closed immersion
$$
\mathbf{P}(\mathcal{E}) \longrightarrow
\mathbf{P}(\mathcal{O}_S^{\oplus n}) = \mathbf{P}^{n - 1}_S
$$
as desired. In general, let $\mathcal{L}$ be an ample invertible sheaf on $S$.
By Properties, Proposition \ref{properties-proposition-characterize-ample}
there exists an integer $n$ such that
$\mathcal{E} \otimes_{\mathcal{O}_S} \mathcal{L}^{\otimes n}$
is globally generated by finitely many sections. Since
$\mathbf{P}(\mathcal{E}) =
\mathbf{P}(\mathcal{E} \otimes_{\mathcal{O}_S} \mathcal{L}^{\otimes n})$ by
Constructions, Lemma \ref{constructions-lemma-twisting-and-proj}
this finishes the proof.
\end{proof}

\begin{lemma}
\label{lemma-proper-ample-is-proj}
Let $f : X \to S$ be a universally closed morphism.
Let $\mathcal{L}$ be an $f$-ample invertible $\mathcal{O}_X$-module.
Then the canonical morphism
$$
r : X
\longrightarrow
\underline{\text{Proj}}_S
\left(
\bigoplus\nolimits_{d \geq 0} f_*\mathcal{L}^{\otimes d}
\right)
$$
of Lemma \ref{lemma-characterize-relatively-ample} is an isomorphism.
\end{lemma}

\begin{proof}
Observe that $f$ is quasi-compact because the existence
of an $f$-ample invertible module forces $f$ to be quasi-compact.
By the lemma cited the morphism $r$ is an open immersion.
On the other hand, the image of $r$ is closed by
Lemma \ref{lemma-image-proper-scheme-closed}
(the target of $r$ is separated over $S$ by Constructions,
Lemma \ref{constructions-lemma-relative-proj-separated}).
Finally, the image of $r$ is dense by
Properties, Lemma \ref{properties-lemma-ample-immersion-into-proj}
(here we also use that it was shown in the proof of
Lemma \ref{lemma-characterize-relatively-ample}
that the morphism $r$ over affine opens of $S$
is given by the canonical morphism of
Properties, Lemma \ref{properties-lemma-map-into-proj}).
Thus we conclude that $r$ is a surjective open immersion, i.e.,
an isomorphism.
\end{proof}

\begin{lemma}
\label{lemma-proper-ample-delete-affine}
Let $f : X \to S$ be a universally closed morphism. Let $\mathcal{L}$
be an $f$-ample invertible $\mathcal{O}_X$-module. Let
$s \in \Gamma(X, \mathcal{L})$. Then $X_s \to S$ is an affine morphism.
\end{lemma}

\begin{proof}
The question is local on $S$ (Lemma \ref{lemma-characterize-affine})
hence we may assume $S$ is affine. By Lemma \ref{lemma-proper-ample-is-proj}
we can write $X = \text{Proj}(A)$ where $A$ is a graded
ring and $s$ corresponds to $f \in A_1$ and
$X_s = D_+(f)$ (Properties, Lemma \ref{properties-lemma-map-into-proj})
which proves the lemma by construction of $\text{Proj}(A)$, see
Constructions, Section \ref{constructions-section-proj}.
\end{proof}











